\subsection*{Question 1}
Dans le repère, les coordonnées du point A sont A$(-2 ; 2)$.

\subsection*{Question 2}
Le point O est le milieu du segment [MK] donc le symétrique du point K par la symétrie centrale de centre O est M.

\subsection*{Question 3}
La formule pour calculer l'aire d'un triangle est \quad $\dfrac{\text{Base}\times \text{Hauteur}}{2}$.

Ici, en prenant [AC] comme base, la hauteur correspondante est [BD].

Le calcul attendu est donc :\quad $\mathrm{AC} \times \mathrm{BD}\div 2$.

Bonne réponse : \textbf{B.}

\subsection*{Question 4}
Les résultats obtenus sont : \emph{pile, pile, face, pile, face, face, face, face, pile, face}, soit 4 fois \emph{pile} sur 10 lancers.

La fréquence d'apparition de \og \emph{pile} \fg{} est :\quad  $\dfrac{\text{Nombre de \og{}\emph{pile}\fg{}}}{\text{Nombre de lancers}} = \dfrac{4}{10}$.

Bonne réponse : \textbf{C.}

\emph{Remarque} : la \og bonne \fg{} réponse est en fait $\dfrac25$.

\subsection*{Question 5}

On note $n$ un nombre entier.

%	Parmi les propositions suivantes, quelle expression donne la moitié de $n$ ?
La moitié de $n$ est $\dfrac{n}{2}$.

Bonne réponse : \textbf{D.}

\emph{Remarque :} $2n$ est le double de $n$, $n^2$ est son carré, et $n+2$ est la somme de $n$ et de 2.

\subsection*{Question 6}

1 h vaut 60 min, donc $\dfrac{1}{10}$ d'heure vaut $\dfrac{1}{10}\times 60 = \dfrac{60}{10}$, soit  $\np[min]{6}$.

Bonne réponse : \textbf{C.}

%	\begin{tblr}{*{4}{X}}
%		\textbf{A.}\quad 0,1 min &
%		\textbf{B.}\quad 1 min &
%		\textbf{C.}\quad 6 min &
%		\textbf{D.}\quad 10 min
%	\end{tblr}

\subsection*{Question 7}

La moyenne de cette série est $\dfrac{10+10+12+16}{4} = \dfrac{48}{4} = {12}$

\subsection*{Question 8}

La somme des mesures des angles d'un triangle est égale à $180$ en degré donc \quad
$x+x+120 = 180$ \quad d'où \quad $2x+120 = 180$,

ou en ajoutant $- 120$ à chaque membre :\quad $2x = 60$\quad donc\quad $x = 30$ (en degrés).

\subsection*{Question 9}

Si on note $v$ la vitesse moyenne, $d$, la distance parcourue et $t$ le temps de parcours, on a :

$v = \dfrac{d}{t}$ \quad donc \quad $vt = d$ puis $t = \dfrac{d}{v} = \dfrac{15}{20} = \dfrac{5\times 3}{5\times 4} = {\dfrac{3}{4}} = \dfrac{3 \times 15}{4 \times 15} = \dfrac{45}{60}$ en heure.
\quad
Le trajet dure  donc 45 min.

\needspace{15\baselineskip}
% saut de page si moins de 10 lignes disponibles
